\section{Hyperparameter Search: Run~2}
\label{app:hyperparam}

Run~2 extended the proof-of-concept by introducing all five attack methods
(Loss, Zlib, Ratio, SAMA, XGBoost, MLP) on the GitHub domain and varying the
number of fine-tuning epochs to find the operating point with the best
membership discriminability.

\subsection*{Epoch Ablation}

The central finding is that discriminability is \emph{non-monotone} in training
epochs.  With too few epochs the model has not memorised sufficiently to create
an ELBO gap; with too many epochs (15 as in Run~1), over-memorisation causes
the SAMA baseline to collapse and makes the domain generally ``easy'' regardless
of membership status.  The 5-epoch operating point produced ELBO member gaps
of 2.19--2.75 nats across domains—within SAMA's designed operating range—while
maintaining strong XGBoost discriminability.

\subsection*{Final Hyperparameter Configuration}

\begin{center}
\small
\begin{tabular}{ll}
\toprule
\textbf{Parameter} & \textbf{Value} \\
\midrule
Fine-tuning epochs         & 5 \\
Learning rate              & $10^{-4}$ \\
Weight decay               & 0.01 \\
Batch size                 & 8 \\
Precision                  & bfloat16 \\
Masking ratios $\boldsymbol\alpha$ & \{0.05, 0.20, 0.35, 0.50\} \\
Masks per ratio $K$        & 8 \\
Feature dimensionality     & 46 \\
SAMA subsets $N$           & 128 \\
SAMA steps $T$             & 4 \\
XGBoost trees              & 200 \\
XGBoost max depth          & 4 \\
XGBoost learning rate      & 0.05 \\
MLP hidden layers          & (256, 128, 64) \\
MLP dropout                & 0.3 \\
CV folds                   & 5 \\
Bootstrap samples          & 1000 \\
\bottomrule
\end{tabular}
\end{center}

All experiments in Run~3 and Run~4 use these exact hyperparameters.
